% =========================================================================
% APPENDICE C — CODICE E RIPRODUCIBILITÀ
% =========================================================================

\chapter{Codice e riproducibilità}
\label{app:code}

Questa appendice documenta il codice che sta dietro alla tesi, così che la pipeline analitica e
l'artefatto rappresentazionale possano essere ispezionati e riprodotti. Copre i repository,
l'ambiente computazionale, l'accesso ai dati, i passi per riprodurre l'analisi e un registro di
provenienza che fa risalire ogni elemento percepibile dell'artefatto alla sua fonte nei dati, al
principio di design che lo autorizza e alla sua origine computazionale. L'impegno guida è quello che
percorre l'intera tesi: ogni valore mostrato deriva dai dati, e il percorso dai dati al display è
aperto all'ispezione. Un'affermazione scientifica che non può essere verificata in modo indipendente
non è, in un senso significativo, un'affermazione scientifica.

% -------------------------------------------------------------------------
\section{Repository e struttura del progetto}
\label{app:c-repos}

Il progetto è suddiviso su tre repository pubblici anziché su uno solo. La separazione è deliberata:
la pipeline analitica, l'artefatto rappresentazionale e la tesi scritta hanno cicli di vita di
sviluppo diversi, linguaggi e toolchain diversi, e storie più pulite se tenuti separati. Tutti e tre
sono ospitati pubblicamente sotto \url{https://github.com/enricovaccari}.

\subsection*{Pipeline analitica}
\textit{Repository:} \texttt{amoc-thesis-pipeline} ---
\url{https://github.com/enricovaccari/amoc-thesis-pipeline}.

Questo repository contiene i notebook e i moduli che implementano la fase di \emph{rilevamento}.
L'analisi è organizzata come una sequenza di notebook numerati, ciascuno uno stadio autonomo il cui
output alimenta il successivo:

\begin{itemize}
    \item \texttt{01\_data\_first\_look} --- caricamento e ispezione del record RAPID;
    \item \texttt{02\_rapid\_vs\_copernicus} --- confronto dell'osservazione con l'ensemble di
          rianalisi;
    \item \texttt{03\_PCA\_on\_vertical\_profiles} --- la decomposizione lineare della struttura
          verticale;
    \item \texttt{04\_anomalies\_early\_warning\_signals} --- il rilevamento delle anomalie e
          l'analisi corretta dei segnali di preallarme;
    \item \texttt{05\_autoencoder\_vertical\_profiles} --- il confronto con un modello non lineare
          che verifica la planarità della varietà;
    \item \texttt{06\_export\_for\_translate} --- assemblaggio del bundle di dati consumato
          dall'artefatto.
\end{itemize}

\noindent La funzionalità condivisa è fattorizzata in moduli di supporto --- \texttt{data\_loading.py}
per i reader e \texttt{paths.py} per la disposizione delle directory --- anziché duplicata tra i
notebook. Il notebook finale esporta un unico bundle JSON,
\texttt{amoc\_translate\_bundle.json}, che è l'unica fonte di verità per ogni valore che l'artefatto
mostra (\S\ref{app:c-ledger}).

\subsection*{Artefatto rappresentazionale}
\textit{Repository:} \texttt{amoc-thesis-translate} ---
\url{https://github.com/enricovaccari/amoc-thesis-translate}, distribuito tramite GitHub Pages.

Questo repository contiene l'artefatto basato su browser: le tre condizioni rappresentazionali
costruite per lo studio, animate nel browser, con il suono dalla Web Audio API nativa. La sua
struttura separa le tre scene dalla macchineria condivisa da cui dipendono:

\begin{itemize}
    \item i file di scena, uno statico e uno esperienziale per scena ---
          \nolinkurl{column-static.html}, \nolinkurl{theengine-static.html} e
          \nolinkurl{thethreshold-static.html}, con le loro controparti \texttt{-experiential};
    \item uno strato di moduli condivisi --- \texttt{shared/bundle.js} carica il bundle di dati,
          \texttt{shared/reconstruct.js} ricostruisce ogni profilo attraverso
          l'Equazione~\ref{eq:reconstruction}, e \texttt{shared/intro-*.js} e
          \texttt{shared/context-*.js} portano il testo di inquadramento importato in modo identico
          byte per byte da entrambe le condizioni.
\end{itemize}

\noindent Questa architettura a moduli condivisi è ciò che permette alla tesi di affermare che le
condizioni statica ed esperienziale sono allineate per provenienza, inquadramento e unità di misura:
attingono allo stesso bundle e alla stessa funzione di ricostruzione, e differiscono solo nel modo in
cui lo rendono --- la versione esperienziale aggiunge la dimensione temporale che lo studio esamina.
Quella provenienza allineata è garantita da codice condiviso anziché asserita --- un bundle, una
ricostruzione, import condivisi, non codice duplicato.

\subsection*{Tesi scritta}
\textit{Repository:} \texttt{amoc-thesis-writing} ---
\url{https://github.com/enricovaccari/amoc-thesis-writing}. I sorgenti \LaTeX{} e il database
bibliografico. Non è un contributo di ricerca e non è necessario per riprodurre né l'analisi né
l'artefatto; è elencato per completezza.

% -------------------------------------------------------------------------
\section{Ambiente computazionale}
\label{app:c-environment}

La pipeline analitica gira in Python~3.11 dentro un ambiente Conda chiamato
\texttt{amoc-thesis}, registrato come kernel Jupyter dedicato così che i notebook girino in modo
riproducibile da dentro Visual Studio Code. Le librerie principali sono \texttt{xarray} e
\texttt{netCDF4} per i dati oceanici, \texttt{numpy}, \texttt{pandas}, \texttt{scipy} e
\texttt{statsmodels} per il lavoro numerico e sulle serie temporali, \texttt{scikit-learn} per la PCA
e l'Isolation Forest, \texttt{PyTorch} per l'autoencoder, e \texttt{matplotlib} per i grafici. Le
versioni esatte sono fissate nel file d'ambiente del repository, da cui l'ambiente può essere
ricreato per intero.

Per riprodurre l'analisi, l'ambiente è ricreato da quel file e i notebook numerati sono eseguiti in
ordine. Eseguirli in sequenza rigenera il bundle di dati, da cui l'artefatto può poi essere servito
localmente come file statici.

L'artefatto rappresentazionale non richiede alcun ambiente di build: è scritto in JavaScript puro e
nella Web Audio API nativa, senza dipendenze di rendering esterne, ed è servito o come file statici
in locale o tramite GitHub Pages.

% -------------------------------------------------------------------------
\section{Accesso ai dati}
\label{app:c-data}

Le due fonti di dati sono ottenute dai rispettivi fornitori anziché incluse con il codice, sia per la
dimensione dei file sia perché le loro licenze governano la ridistribuzione. Il record osservativo
RAPID è distribuito pubblicamente dal British Oceanographic Data Centre \citep{moat2026}; l'ensemble
di rianalisi GREP del Copernicus Marine Service è accessibile sotto la licenza dati Copernicus
\citep{copernicus_grep}. Il README della pipeline documenta il percorso di download per ciascuno e
dove ogni file va collocato così che i notebook risolvano i loro path; i dati grezzi stanno fuori dal
controllo di versione per disegno.

% -------------------------------------------------------------------------
\section{Registro di provenienza}
\label{app:c-ledger}

Il registro di provenienza (Tabella~\ref{tab:ledger}) è l'affermazione di riproducibilità resa
concreta. Per ogni elemento percepibile dell'artefatto, nomina il campo del bundle di dati da cui
deriva, il principio della grammatica di design che ne autorizza la codifica percettiva e la sua
origine computazionale --- se empirica, derivata classicamente o prodotta dal machine learning. Nulla
di percepibile è omesso tranne l'impalcatura non informativa (colori, font, tinte), che non porta
dati. Il registro è ciò che trasforma il principio di fedeltà da promessa in qualcosa che un lettore
può verificare riga per riga.

% ============================================================
% Il registro di provenienza.  Richiede \usepackage{tabularx, booktabs}.
% ============================================================
\begin{table}[tbp]
\centering
\caption[Registro di provenienza]{Registro di provenienza per le tre scene. Ogni
elemento percepibile, la sua fonte nel bundle di dati, il principio di design che
lo autorizza e la sua origine computazionale. Colori, font e tinte sono
impalcatura non informativa e sono omessi.}
\label{tab:ledger}
\footnotesize
\renewcommand{\arraystretch}{1.3}
\begin{tabularx}{\textwidth}{@{}p{1.2cm} X p{3.4cm} p{1.1cm} X@{}}
\toprule
\textbf{Scena} & \textbf{Elemento} & \textbf{Campo bundle} & \textbf{Prin.} &
\textbf{Origine} \\
\midrule
I   & Curva $\Psi(z)$, statica (C1) & \texttt{mean\_profile} via Eq.~\ref{eq:reconstruction}, $\mathrm{pc}=0$ & P1 & Empirica (media del record) \\
I   & Curva $\Psi(z)$, animata (C2/C3) & \texttt{mean\_profile} + \texttt{trajectory\_monthly} & P1, P2 & ML (traiettoria latente; H2) \\
I   & Campo di particelle & $\partial\Psi/\partial z$ del profilo corrente (Eq.~\ref{eq:velocity}) & P2 & Classica (derivata) \\
I   & Texture dell'anomalia & \texttt{anomalies.clusters} (score) & P3 & ML (Isolation Forest) \\
I   & Suono & \texttt{trajectory\_monthly.pc1} & P5 & Input ML, sintesi classica \\
I   & Lettura profondità / picco & \texttt{depth\_m}; max del profilo & P1 & Derivata \\
\addlinespace
II  & Punti mensili & \texttt{trajectory\_monthly.pc1/pc2} & P1 & ML (coordinate latenti) \\
II  & Moto del percorso & ordine di \texttt{trajectory\_monthly} & P2 & ML (traiettoria latente) \\
II  & Nube sub-mensile & \texttt{manifold\_cloud.pc} & P1 & ML (coordinate latenti) \\
II  & Agitazione dell'anomalia & \texttt{anomalies.clusters} (score) & P3 & ML (Isolation Forest) \\
II  & Bagliore di equilibrio & centroide di \texttt{manifold\_cloud} & P1 & Derivata (media) \\
II  & Suono & \texttt{trajectory\_monthly.pc1} & P5 & Input ML, sintesi classica \\
\addlinespace
III & Tre linee dei membri & \texttt{copernicus\_ensemble.members} & P1 & Rianalisi esterne \\
III & Mezzo vivente & \texttt{upper} $-$ \texttt{lower} per mese & P4 & Classica (dispersione) \\
III & Temperatura di colore & \texttt{upper} $-$ \texttt{lower} per mese & P4 & Classica (dispersione) \\
III & Viste di aggregazione & \texttt{members} (media mobile, inviluppo) & P1 & Derivata \\
III & Texture sonora & \texttt{upper} $-$ \texttt{lower} per mese & P5 & Classica, sintesi classica \\
\bottomrule
\end{tabularx}
\end{table}

\noindent La sonificazione delle Scene~I e~II prende lo stesso input, $\mathrm{pc}_1$, e
differisce solo nella sintesi: la Scena~I lo mappa su un impulso lento, la Scena~II su un tono
continuo che segue il marcatore. Entrambe sono disattivate di default.

% -------------------------------------------------------------------------
\section{Strumenti e workflow}
\label{app:c-tools}

Oltre al codice in sé, la tesi è stata prodotta con un piccolo insieme di strumenti, registrati qui
per completezza e riproducibilità.

La scrittura è stata fatta in Visual Studio Code con l'estensione LaTeX Workshop, compilando con
\TeX{}~Live~2026 attraverso \texttt{latexmk}, con \texttt{biblatex} e \texttt{biber} per la
bibliografia. I riferimenti sono stati gestiti in Zotero con l'estensione Better BibTeX, che esporta
un unico \texttt{references.bib} dalla collezione della tesi usando una convenzione di chiavi di
citazione coerente; l'estensione è connessa direttamente all'editor così che le citazioni fluiscano
nei sorgenti \LaTeX{} senza copiature manuali. Il controllo di versione lungo tutto il lavoro ha
usato Git con GitHub, un repository per componente come descritto sopra.

Lo studio empirico si è affidato a Google Forms per le quattro varianti del questionario e a
Microsoft Excel per il dataset delle risposte codificate. L'analisi è girata in notebook Jupyter
dentro l'ambiente Conda. L'organizzazione del progetto e la presa di appunti sono state tenute in
Notion, e la comunicazione con il relatore è avvenuta su Slack. Dove le ricerche bibliografiche
avevano bisogno di un punto di partenza, Perplexity è stato usato per individuare fonti candidate, poi
verificate e lette di prima mano prima della citazione.

Il ruolo degli assistenti di IA in questo processo, e i confini posti attorno a esso, è discusso
separatamente come questione di etica della ricerca nell'Appendice~\ref{app:ai-use}.
